% 表紙情報はこのファイルで編集する．年度・所属・氏名等は未確定．
\newcommand{\ThesisTitle}{獣害対策トレイルカメラシステムに関する研究}
% 表紙だけで使用する改行位置．題目変更時は上の命令と併せて編集する．
\newcommand{\ThesisCoverTitle}{獣害対策トレイルカメラシステム\\に関する研究}
\newcommand{\ThesisYear}{［年度］}
\newcommand{\ThesisDate}{［提出年月］}
\newcommand{\ThesisUniversity}{［大学名］}
\newcommand{\ThesisGraduateSchool}{［研究科・専攻名］}
\newcommand{\ThesisField}{［分野・コース名］}
\newcommand{\ThesisUniversityEnglish}{[University name]}
\newcommand{\ThesisSchoolEnglish}{[Graduate school name]}
\newcommand{\ThesisSupervisor}{［指導教員氏名・職位］}
\newcommand{\ThesisStudentID}{［学籍番号］}
\newcommand{\ThesisAuthor}{［氏名］}

% テンプレートの表紙用コマンド．年度には「年度」を付けず記入する．
\title{\shortstack{\ThesisCoverTitle}}
\author{\ThesisAuthor}
\authorintitle{\ThesisAuthor}
\date{\ThesisDate}
\nendo{\ThesisYear}
\studentid{\ThesisStudentID}
\advisor{\ThesisSupervisor}
\affiliation{\shortstack{\ThesisUniversity\quad\ThesisGraduateSchool\\\ThesisField}}
\pdfstringdefDisableCommands{\def\shortstack#1{#1}\def\\{ }}
